\documentclass[../../main.tex]{subfiles}
\begin{document}

\begin{flushright}
\footnotesize\textit{【自動文字起こし・要確認】}
\end{flushright}

[解] $v = const$ とする。$100 \text{ km}$ 走るのに要する時間は $\dfrac{100}{v} \ [h]$ である。 $\cdots \text{①}$ \\
はじめにガソリンが $x_0 \ [l]$ あったとして、題意から
\[ \dfrac{dx}{dt} = - \dfrac{100 + x}{100} \cdot e^{kv} \]
\[ \dfrac{1}{100 + x} dx = - \dfrac{e^{kv}}{100} dt \]
積分して、初期条件から、
\[ x + 100 = (x_0 + 100)e^{-\dfrac{e^{kv}}{100}t} \cdots \text{②} \]
したがって、①から、総ガソリン消費量 $Y$ は、②より
\[ Y = x_0 - x|_{t=\dfrac{100}{v}} \]
\[ = x_0 - \left\{ (x_0+100) e^{-\dfrac{e^{kv}}{v}} - 100 \right\} \quad \left( A = \dfrac{e^{kv}}{v} \text{ とおく} \right) \]
\[ = (1 - e^{-A})x_0 - 100(-1 + e^{-A}) \]

$-\dfrac{e^{kv}}{v} < 0$ から $Y(x_0)$ は 1次係数正の $x_0$ の1次式 $\cdots \text{④}$ \\
又、$t = \dfrac{100}{v}$ で $x \ge 0$ だから、②より
\[ (x_0 + 100) e^{-\dfrac{e^{kv}}{v}} \ge 100 \]
\[ x_0 \ge 100(e^{\dfrac{e^{kv}}{v}} - 1) \equiv P \cdots \text{⑤} \]
④、⑤より、$Y(v)$ を $\min$ にする $x_0$ は $P$ である。
\[ \min Y(v) = (1 - e^{-A})100(e^A - 1) - 100(-1 + e^{-A}) \]
\[ = 100[e^A - 1] \cdots \text{⑥} \]

$e^x$ は単調増加だから $A$ が $\min$ の時 $Y$ も $\min$ である
\[ \dfrac{dA}{dv} = \dfrac{k e^{kv} v - e^{kv}}{v^2} = \dfrac{kv - 1}{v^2} e^{kv} \]
から、$v = \dfrac{1}{k} \ (>0)$ で $A$ は $\min$ である。⑤からこの時
\[ P = 100(e^{ek} - 1) \] \hfill(終)

\end{document}
